% Generated draft: requires mathematical and source review.
\documentclass[11pt]{article}
\usepackage[margin=25mm]{geometry}
\usepackage{fontspec}
\setmainfont{Latin Modern Roman}
\usepackage{amsmath,amssymb,mathtools,unicode-math}
\setmathfont{Latin Modern Math}
\usepackage{enumitem,needspace,xurl,hyperref}
\setlength{\parindent}{0pt}
\setlength{\parskip}{5pt}
\setlength{\emergencystretch}{3em}
\newcommand{\N}{\mathbb{N}}
\newcommand{\Z}{\mathbb{Z}}
\newcommand{\Q}{\mathbb{Q}}
\newcommand{\R}{\mathbb{R}}
\newcommand{\C}{\mathbb{C}}
\newcommand{\F}{\mathbb{F}}
\newcommand{\A}{\mathbb{A}}
\newcommand{\PP}{\mathbb{P}}
\newcommand{\E}{\mathbb{E}}
\newcommand{\eps}{\varepsilon}
\newcommand{\dd}{\,\mathrm{d}}
\begin{document}
\section*{UnsolvedMath ID: 30001344}
\textbf{Source problem code:} catalogue entry
\subsection{Definitions and mathematical statement}
\emph{Terminology and scope.} The supplied record concerns differential equations with irregular singularities and uses the expression ``wild geometry.'' It supplies no definition of wild geometry, no class of differential equations, no underlying spaces or ground field, and no criterion for usefulness in algebraic geometry. These choices therefore remain unspecified.

\emph{Source question.} ``What is the usefulness for algebraic geometry of wild geometry for differential equations with irregular singularities?''

This is a research-direction question. The supplied record does not specify a mathematical assertion with definite hypotheses, quantifiers, and conclusion. A precise theorem or conjecture cannot be extracted without adding information: in particular, the intended geometric constructions, the algebraic-geometric targets, and the properties to be established must be identified. No such choices are imposed in this draft.
\subsection{Short English statement}
The question asks what applications wild geometry associated with differential equations with irregular singularities has in algebraic geometry. The intended objects and meaning of an application are unspecified, so the record supports a broad research question requiring further formulation.
\subsection{Sources}
\begin{itemize}
\item Local problem registry: \texttt{data/unsolvedmath/problems.json}, permanent UnsolvedMath ID \(30001344\), source code \texttt{catalogue entry}; title: ``Applications of Wild Geometry to Algebraic Geometry.'' This draft uses the supplied complete registry record.
\item Supplied report citation: \emph{Singularities (2010)}, Oberwolfach Reports, DOI \texttt{10.4171/owr/2009/43}, pp.~2405--2470. Supplied source URL and literature-source link, labeled ``original Oberwolfach report'': \url{https://doi.org/10.4171/owr/2009/43}. The record describes Oberwolfach Reports as documenting mathematical workshops and their open-problem discussions.
\item Attribution: the record supplies proposed year \(2009\) but no named proposer. Its source citation uses \(2010\), whereas its DOI and year tag contain \(2009\); the supplied evidence does not explain this dating distinction.
\item Supplied literature assessment, checked \textbf{2026-08-21}: the record labels the question open and states that irregular Riemann--Hilbert and wild Hodge theories provide substantial applications, while the extracted broad usefulness question is not a single settled theorem. This is supplied dated curation derived from the OWR record, not an independently verified literature assessment. The record explicitly calls for independent verification of time-sensitive status and cited claims.
\item The registry marks the statement as exact and reports that manual review found the extraction self-contained, mathematically coherent, and free of LaTeX or report-layout artifacts. These are supplied extraction-verification claims; they do not resolve the missing mathematical specifications.
\end{itemize}
\end{document}
